\section*{Capitulo 9. Derivadas y aplicaciones}

\begin{shortanswer}{[GX-C09.S01-01]}
\[
g'(x)=\frac{3x^2+6x+1}{(x+1)^2}.
\]
\end{shortanswer}

\begin{shortanswer}{[GX-C09.S01-02]}
\[
h'(x)=\frac{1}{\sqrt{2x+5}}.
\]
\end{shortanswer}

\begin{shortanswer}{C09.S01 practica autonoma}
\begin{enumerate}[label=\textbf{PX-C09.S01-\arabic*:}, leftmargin=*]
  \item \(4x^3-3\)
  \item \(3x^2-6x+1\)
  \item \(-\dfrac{5}{(x-2)^2}\)
  \item \(\dfrac{x}{\sqrt{x^2+4}}\)
  \item \(3e^{3x}\)
  \item \(\dfrac{2x}{x^2+1}\)
  \item \(\cos^2 x-\sen^2 x\)
  \item \(6x(x^2-1)^2\)
\end{enumerate}
\end{shortanswer}

\begin{shortanswer}{[GX-C09.S02-01]}
Decrece en \((-\infty,2)\), crece en \((2,\infty)\) y tiene minimo en \((2,-3)\).
\end{shortanswer}

\begin{shortanswer}{[GX-C09.S02-02]}
Crece en todo \(\mathbb{R}\), es concava hacia abajo en \((-\infty,0)\), hacia arriba en
\((0,\infty)\) y tiene punto de inflexion en \((0,0)\).
\end{shortanswer}

\begin{shortanswer}{C09.S02 practica autonoma}
\begin{enumerate}[label=\textbf{PX-C09.S02-\arabic*:}, leftmargin=*]
  \item Crece en \((-\infty,2)\), decrece en \((2,\infty)\) y tiene maximo en \((2,4)\)
  \item Crece en \((-\infty,0)\cup(4,\infty)\), decrece en \((0,4)\), maximo en \((0,0)\),
  minimo en \((4,-32)\) e inflexion en \((2,-16)\)
  \item Decrece en \((-\infty,-1)\cup(0,1)\), crece en \((-1,0)\cup(1,\infty)\), maximo en
  \((0,0)\) y minimos en \((-1,-1)\) y \((1,-1)\)
  \item Maximo absoluto \(5\) en \(x=0\); minimo absoluto \(1\) en \(x=2\)
  \item Crece en \((0,\infty)\), es concava hacia abajo en \((0,\infty)\) y no tiene inflexion
  \item Decrece en \((-\infty,0)\), crece en \((0,\infty)\) y tiene minimo en \((0,1)\)
\end{enumerate}
\end{shortanswer}

\begin{shortanswer}{[GX-C09.S03-01]}
Tangente: \(y=-2x\). Normal: \(y=\dfrac{x}{2}+\dfrac{5}{2}\).
\end{shortanswer}

\begin{shortanswer}{[GX-C09.S03-02]}
Tangente: \(y=-2\). Normal: \(x=1\).
\end{shortanswer}

\begin{shortanswer}{C09.S03 practica autonoma}
\begin{enumerate}[label=\textbf{PX-C09.S03-\arabic*:}, leftmargin=*]
  \item Tangente: \(y=6x-4\). Normal: \(y=-\dfrac{x}{6}+\dfrac{25}{3}\)
  \item Tangente: \(y=-\dfrac{x}{2}+\dfrac{7}{2}\). Normal: \(y=2x-4\)
  \item Tangente: \(y=\dfrac{x}{4}+\dfrac{5}{4}\). Normal: \(y=-4x+14\)
  \item Tangente: \(y=x+1\). Normal: \(y=-x+1\)
  \item Tangente: \(y=x-1\). Normal: \(y=-x+1\)
  \item Tangente: \(y=0\). Normal: \(x=0\)
\end{enumerate}
\end{shortanswer}

\begin{shortanswer}{[GX-C09.S04-01]}
\[
y=2x-9.
\]
\end{shortanswer}

\begin{shortanswer}{[GX-C09.S04-02]}
\[
y=\frac{x}{4}+2.
\]
\end{shortanswer}

\begin{shortanswer}{C09.S04 practica autonoma}
\begin{enumerate}[label=\textbf{PX-C09.S04-\arabic*:}, leftmargin=*]
  \item \(y=4x-3\)
  \item \(y=12x+16\) y \(y=12x-16\)
  \item \(y=x-1\)
  \item \(y=e^2x-e^2\)
  \item \(y=2\) y \(y=-2\)
  \item \(a=3\) y la tangente es \(y=5x-1\)
\end{enumerate}
\end{shortanswer}

\begin{shortanswer}{[GX-C09.S05-01]}
\[
a=3.
\]
\end{shortanswer}

\begin{shortanswer}{[GX-C09.S05-02]}
\[
b=\frac{3}{2},
\]
y en \(x=-1\) aparece un maximo local.
\end{shortanswer}

\begin{shortanswer}{C09.S05 practica autonoma}
\begin{enumerate}[label=\textbf{PX-C09.S05-\arabic*:}, leftmargin=*]
  \item \(a=-3\)
  \item \(b=-4\)
  \item \(c=-3\)
  \item \(k=2\)
  \item \(m=-3\)
  \item \(p=2\)
\end{enumerate}
\end{shortanswer}

\begin{shortanswer}{[GX-C09.S06-01]}
Es continua en \(x=0\), pero no es derivable.
\end{shortanswer}

\begin{shortanswer}{[GX-C09.S06-02]}
Es continua para todo \(k\) y es derivable solo si \(k=1\).
\end{shortanswer}

\begin{shortanswer}{C09.S06 practica autonoma}
\begin{enumerate}[label=\textbf{PX-C09.S06-\arabic*:}, leftmargin=*]
  \item Es continua y derivable
  \item Es continua, pero no derivable
  \item \(a=0\)
  \item Es continua, pero no derivable
  \item \(b=-4\), y con ese valor es continua y derivable
  \item \(m=2\)
\end{enumerate}
\end{shortanswer}

\begin{shortanswer}{[GX-C09.S07-01]}
Dos lados de \(4{,}5\) m y el lado paralelo al muro de \(9\) m. Area maxima:
\(40{,}5\ \text{m}^2\).
\end{shortanswer}

\begin{shortanswer}{[GX-C09.S07-02]}
El minimo se alcanza en \(x=3\) y vale \(6\).
\end{shortanswer}

\begin{shortanswer}{C09.S07 practica autonoma}
\begin{enumerate}[label=\textbf{PX-C09.S07-\arabic*:}, leftmargin=*]
  \item El cuadrado de lado \(6\). Area maxima: \(36\)
  \item Se alcanza con \(5\) y \(5\). Producto maximo: \(25\)
  \item El minimo es \(4\) y se alcanza en \(x=2\)
  \item Lados perpendiculares de \(7{,}5\) m y lado paralelo de \(15\) m. Area maxima:
  \(\dfrac{225}{2}\ \text{m}^2\)
  \item Maximo \(36\) en \(x=6\)
  \item \(x=\dfrac{4}{3}\) cm
\end{enumerate}
\end{shortanswer}

\begin{shortanswer}{[GX-C09.S08-01]}
Crece en \((-\infty,0)\cup(2,\infty)\), decrece en \((0,2)\), tiene maximo en \((0,0)\),
minimo en \((2,-4)\) e inflexion en \((1,-2)\).
\end{shortanswer}

\begin{shortanswer}{[GX-C09.S08-02]}
Se obtienen en \(x=-1\) y \(x=1\). Las tangentes son \(y=4\) y \(y=0\).
\end{shortanswer}

\begin{shortanswer}{C09.S08 practica autonoma}
\begin{enumerate}[label=\textbf{PX-C09.S08-\arabic*:}, leftmargin=*]
  \item Crece en \((-\infty,-\sqrt{2})\cup(\sqrt{2},\infty)\) y decrece en
  \((-\sqrt{2},\sqrt{2})\). Tiene maximo en \((-\sqrt{2},4\sqrt{2})\),
  minimo en \((\sqrt{2},-4\sqrt{2})\) e inflexion en \((0,0)\)
  \item Tangente: \(y=-3x+1\). Normal: \(y=\dfrac{x}{3}+1\)
  \item En \(x=-2\) y \(x=2\). Las tangentes son \(y=16\) y \(y=-16\)
  \item Maximo absoluto \(2\) en \(x=-1\) y \(x=2\); minimo absoluto \(-2\) en \(x=-2\) y
  \(x=1\)
  \item Puntos criticos en \(x=-\sqrt{2},0,\sqrt{2}\). Maximo local en \((0,0)\) y minimos
  locales en \((-\sqrt{2},-4)\) y \((\sqrt{2},-4)\)
  \item \(a=2,\ b=-1\)
\end{enumerate}
\end{shortanswer}

\begin{shortanswer}{[CH-C09-01]}
La tangente es horizontal en \(t=3\). Su ecuacion es \(y=14\). En ese instante la pelota
alcanza una altura maxima instantanea y la velocidad vertical es cero.
\end{shortanswer}
